% sections/03_representation_routes.tex  —  Representation Routes
% \input{} into main.tex; do NOT wrap in \begin{document}.

\section{Representation Routes}
\label{sec:repr}

The raw measurement—a pair of real-valued vectors (magnitude and phase)
over \num{4001} frequency points—must be transformed into a feature
representation before it can be consumed by a learning algorithm.  The
choice of representation is not neutral: it determines which spectral
information is preserved, which is discarded, and which artefacts are
introduced.  This section develops two complementary representation
families, designated \emph{Route~A} (1D spectral) and \emph{Route~B}
(2D image-like), and evaluates their trade-offs systematically.

% =====================================================================
\subsection{Route A: 1D Spectral Representations}
\label{sec:repr:routeA}

Route~A treats each $S_{11}$ frequency sweep as a one-dimensional
feature vector whose elements correspond to uniformly spaced frequency
bins.  The spectral ordering is preserved, and no spatial rasterisation
is performed.

The raw data is acquired in native polar format ($|S_{11}|$ in~dB,
$\angle S_{11}$ in degrees).  A purely Cartesian representation
(Real/Imaginary only) would compress the deep \SI{-62}{\decibel}
unloaded resonance feature into a linear amplitude of
$\sim 0.0007$, rendering it indistinguishable from noise for
gradient-based models.  Conversely, retaining the dB magnitude alone
would sacrifice the mathematically continuous phase information needed
by convolutional kernels.  Therefore, Route~A utilises a
\textbf{3-channel hybrid approach}:

% ---------------------------------------------------------------------
\subsubsection{A1: Three-Channel Hybrid Tensor}
\label{sec:repr:routeA:hybrid}

The primary 1D encoding constructs a three-channel tensor from the
native dB/Phase data:
\begin{equation}
    \mathbf{x}_{\text{3ch}} =
    \begin{bmatrix}
        |S_{11}(f_1)|_{\text{dB}}, & \ldots, & |S_{11}(f_{4001})|_{\text{dB}} \\
        \text{Re}[S_{11}(f_1)],    & \ldots, & \text{Re}[S_{11}(f_{4001})] \\
        \text{Im}[S_{11}(f_1)],    & \ldots, & \text{Im}[S_{11}(f_{4001})]
    \end{bmatrix}
    \;\in\;\mathbb{R}^{3 \times 4001},
    \label{eq:three_channel}
\end{equation}
where the Cartesian components are derived from the native polar values:
\begin{align}
    \text{Re}[S_{11}(f_i)] &= 10^{\,|S_{11}(f_i)|_{\text{dB}}\,/\,20}
        \;\cos\!\bigl(\angle S_{11}(f_i) \cdot \pi/180\bigr),
        \label{eq:polar_to_re} \\
    \text{Im}[S_{11}(f_i)] &= 10^{\,|S_{11}(f_i)|_{\text{dB}}\,/\,20}
        \;\sin\!\bigl(\angle S_{11}(f_i) \cdot \pi/180\bigr).
        \label{eq:polar_to_im}
\end{align}

This design is engineered to maximise the information available to
learning algorithms:
\begin{itemize}
    \item \textbf{Channel~1 (dB)} preserves the full \SI{45}{\decibel}
        dynamic range needed to detect the unloaded patch resonance,
        which appears as a massive, high-contrast dip.  In purely linear
        Cartesian coordinates, this \SI{-62}{\decibel} feature would
        collapse to an amplitude of $\sim 0.0007$, making it
        indistinguishable from numerical noise.
    \item \textbf{Channels~2 and~3 (Re, Im)} provide the mathematical
        continuity needed for 1D~CNN kernels without
        $\pm\SI{360}{\degree}$ discontinuities.  Phase wrapping in the
        native polar format creates artificial edges that would trigger
        false feature detections in convolutional layers.
\end{itemize}

This $3 \times 4001$ tensor is the natural input format for
one-dimensional convolutional architectures (1D~CNN, ResNet-1D).
Because convolution is a linear operator, sliding a filter across the
three-channel input allows the network to simultaneously exploit
high-contrast magnitude features and smooth phase-derived signals.

% ---------------------------------------------------------------------
\subsubsection{A2: Flattened Three-Channel Vector}
\label{sec:repr:routeA:flat}

For classical machine learning algorithms (e.g., Logistic Regression,
Random Forest, Gradient Boosted Trees) that expect a one-dimensional
feature vector, the three-channel tensor is flattened:
\begin{equation}
    \mathbf{x}_{\text{flat}} =
    \bigl[\,|S_{11}(f_1)|_{\text{dB}}, \ldots, |S_{11}(f_{4001})|_{\text{dB}},\;
    \text{Re}(f_1), \ldots, \text{Re}(f_{4001}),\;
    \text{Im}(f_1), \ldots, \text{Im}(f_{4001})\,\bigr]^{\!\top}
    \;\in\;\mathbb{R}^{12003}.
    \label{eq:flat_vector}
\end{equation}

This encoding provides exactly the same information as~A1, but the
spatial co-location of the three channels at the same frequency is
lost to the algorithm (except for tree-based models which can route
based on independent features).  The high dimensionality ($d = 12003$)
relative to the sample size ($N = 1603$) necessitates strong
regularisation to prevent overfitting in linear models.

\paragraph{Ablation variants.}
To quantify the contribution of each channel, baseline experiments
should include:
\begin{itemize}
    \item $C = 1$: magnitude only ($|S_{11}|$ in dB, $d = 4001$).
    \item $C = 2$: Cartesian only (Re, Im; $d = 8002$).
    \item $C = 3$: full hybrid (dB, Re, Im; $d = 12003$ flat or
        $3 \times 4001$ tensor).
\end{itemize}

% ---------------------------------------------------------------------
\subsubsection{A3: Derivative Augmented Cartesian Vector}
\label{sec:repr:routeA:deriv}

Spectral derivatives highlight \emph{local} shape features—slopes,
inflection points, and resonance edges—that may be more discriminative
than the absolute amplitude at any single frequency. The first-order finite
difference approximates the spectral derivative for both components:
\begin{align}
    \Delta_k^{\text{Re}} &= \frac{\text{Re}[S_{11}(f_{k+1})] - \text{Re}[S_{11}(f_k)]}{\Delta f}, \label{eq:first_diff_re} \\
    \Delta_k^{\text{Im}} &= \frac{\text{Im}[S_{11}(f_{k+1})] - \text{Im}[S_{11}(f_k)]}{\Delta f}. \label{eq:first_diff_im}
\end{align}

These differences are appended to the flat vector or concatenated as
additional channels in the tensor representation. While first differences
are invariant to additive baseline shifts (providing robustness against
calibration drift), they amplify high-frequency instrumental noise.

% ---------------------------------------------------------------------
\subsubsection{Route A Encoding Summary}
\label{sec:repr:routeA:summary}

\Cref{tab:routeA_summary} compares the Route~A encoding variants.

\begin{table}[htbp]
    \centering
    \caption{%
        Comparison of Route~A (1D spectral) encoding variants.%
    }
    \label{tab:routeA_summary}
    \small
    \begin{tabular}{llcl}
        \toprule
        \textbf{ID} & \textbf{Encoding} &
        \textbf{Dim.} & \textbf{Model Compatibility} \\
        \midrule
        A1 & Three-Channel Hybrid (dB, Re, Im) &
            $3 \times 4001$ & 1D CNN, ResNet-1D \\
        A2 & Flattened Three-Channel &
            $12003$ & Logistic Regression, SVM, GBT, RF, MLP \\
        A3 & Derivative Augmented &
            $18004$ & All (with higher overfitting risk) \\
        \bottomrule
    \end{tabular}
\end{table}

% =====================================================================
\subsection{Route B: 2D Image-Like Representations}
\label{sec:repr:routeB}

Route~B converts each spectral trace into a two-dimensional image,
enabling the use of established convolutional neural network
architectures originally designed for visual recognition tasks.  The
conversion introduces a \emph{rendering step} that is absent in Route~A,
and the properties of that rendering step materially affect downstream
classification.

% ---------------------------------------------------------------------
\subsubsection{B1: Grayscale Curve Images}
\label{sec:repr:routeB:grayscale}

The most direct conversion plots one of the Cartesian traces (e.g. $\text{Re}[S_{11}(f)]$)
as a curve against frequency on a white background, then rasterises the
plot to a fixed pixel resolution.  Typical target resolutions
are $128\times128$ or $224\times224$ pixels, the latter matching the
default input size of many pretrained image classifiers (e.g.\ ResNet,
VGG, MobileNet).

The resulting single-channel (grayscale) image encodes:
\begin{itemize}
    \item The \emph{vertical position} of the curve at each horizontal
        pixel $\Leftrightarrow$ the amplitude of the chosen component at the corresponding
        frequency.
    \item The \emph{shape} of the curve $\Leftrightarrow$ local spectral
        features (dips, slopes, inflections).
\end{itemize}

\noindent
However, the rendering process introduces several artefacts that have no
physical origin in the measurement:

\begin{itemize}
    \item \textbf{Line thickness and anti-aliasing.}  The width of the
        plotted curve in pixels, and the sub-pixel smoothing applied by
        the rendering engine, create intensity gradients that a CNN may
        learn to exploit.  These gradients are properties of the
        rendering configuration, not of the sensor response.
    \item \textbf{Axis scaling.}  The choice of amplitude axis range
        (e.g.\ $[-1.0,\,1.0]$ vs.\
        $[-0.2,\,0.2]$) affects how much
        of the image canvas the curve occupies.  A tight axis range
        amplifies small variations; a wide range compresses them into
        fewer pixels.
    \item \textbf{Background and margins.}  Axis labels, tick marks,
        grid lines, and border padding consume pixels that carry no
        measurement information but may correlate with class if rendering
        parameters are inadvertently class-dependent.
\end{itemize}

\noindent
To mitigate these risks, the rendering pipeline must use
\emph{identical} axis ranges, line styles, and canvas dimensions for
every sample, and should strip all non-data elements (axes, labels,
borders) from the final image.

% ---------------------------------------------------------------------
\subsubsection{B2: Dual-Channel Rasterised Maps}
\label{sec:repr:routeB:dualchannel}

Analogous to encoding~A1, one may render Real and Imaginary components as two
separate image channels, producing a $H \times W \times 2$ tensor.
Each channel is a grayscale curve image as described in~\Cref{sec:repr:routeB:grayscale},
plotted with independently normalised axis ranges.

This retains the full complex information in image form, and because the Cartesian
components are continuous, it neatly avoids the sharp vertical jumps that would
otherwise be caused by wrapped-phase discontinuities in polar representations.

% ---------------------------------------------------------------------
\subsubsection{B3: Spectrogram-Like and Heatmap Constructions}
\label{sec:repr:routeB:heatmap}

A more creative approach constructs a 2D image not by plotting a curve
but by \emph{tiling} or \emph{reshaping} the spectral data.  Several
strategies are conceivable:

\begin{enumerate}
    \item \textbf{Frequency-axis interpolation.}  The \num{4001}-point
        Cartesian vector is reshaped into a 2D matrix (e.g.\
        $63\times63$ plus padding, or $50\times80$) and rendered as a
        heatmap.  The row and column indices are arbitrary partitions of
        the frequency axis, and the pixel intensity encodes amplitude.
    \item \textbf{Multi-representation tiling.}  Real, Imaginary, and
        derivative traces are rendered as separate horizontal strips and
        stacked vertically into a single image, creating a composite
        ``spectrogram-like'' representation.
    \item \textbf{Recurrence-plot or Gramian-field encoding.}  The
        spectral vector is transformed into a square matrix via a
        pairwise operation (e.g.\ Gramian Angular Field), producing a
        texture image that encodes temporal/spectral correlations.
\end{enumerate}

These approaches can expose inter-frequency relationships that a 1D
convolution might miss, but they introduce a significant design burden:
the reshaping geometry, colour mapping, and interpolation method all
become hyperparameters that are \emph{ad hoc} with respect to the
measurement physics.  Unless a clear empirical advantage is demonstrated,
the additional complexity is difficult to justify for a
three-class problem with \num{1603} samples.

% ---------------------------------------------------------------------
\subsubsection{Rendering Artefact Warning}
\label{sec:repr:routeB:warning}

\begin{quote}
\textbf{Warning.}  Every Route~B representation injects information
into the model's input that is not present in the original measurement
physics.  Line thickness, anti-aliasing kernels, colour-map transfer
functions, axis scaling, and pixel resolution are all rendering
decisions, not sensor observations.  A CNN trained on Route~B images may
learn features that correlate with class labels only because of
consistent rendering choices rather than genuine spectral differences.
This risk is mitigated by strict rendering standardisation but can never
be fully eliminated.  Route~B results should always be compared against
Route~A baselines to assess whether the image pathway adds genuine
discriminative value or merely introduces a confound.
\end{quote}

% =====================================================================
\subsection{Route A vs.\ Route B: Comparative Summary}
\label{sec:repr:comparison}

\Cref{tab:route_comparison} provides a structured comparison of the two
representation families across four evaluation axes.

\begin{table}[htbp]
    \centering
    \caption{%
        Comparison of Route~A (1D spectral) and Route~B (2D image-like)
        representation families.%
    }
    \label{tab:route_comparison}
    \begin{tabular}{p{3.0cm}p{5.2cm}p{5.2cm}}
        \toprule
        \textbf{Axis} & \textbf{Route A (1D Spectral)} &
            \textbf{Route B (2D Image-Like)} \\
        \midrule
        \textbf{Information fidelity} &
            Lossless.  Every frequency-bin value is preserved at full
            floating-point precision. &
            Lossy.  Rasterisation maps continuous-valued amplitudes to
            discrete pixel intensities, and spatial quantisation limits
            the recoverable frequency resolution. \\[6pt]
        \textbf{Computational cost} &
            Low.  No rendering step required.  Preprocessing is limited
            to normalisation and optional feature engineering. &
            Moderate to high.  Rendering each trace to an image adds
            I/O and CPU time.  Larger image resolutions
            ($224\times224$) increase both storage and training cost. \\[6pt]
        \textbf{Model compatibility} &
            Compatible with all classical models (flat vector) and 1D
            neural architectures (1D CNN, ResNet-1D, MLP, KAN).
            Incompatible with 2D CNNs without reshaping. &
            Compatible with 2D CNNs and pretrained image models.
            Incompatible with classical models unless flattened (at
            extreme dimensionality, e.g.\ $224\times224 = 50\,176$). \\[6pt]
        \textbf{Artefact risk} &
            Negligible.  No rendering step means no rendering artefacts.
            Preprocessing artefacts are explicit and auditable. &
            Significant.  Line thickness, anti-aliasing, axis scaling,
            colour mapping, and background elements are all potential
            confounds.  Requires strict standardisation and ablation. \\
        \bottomrule
    \end{tabular}
\end{table}

\subsection{Representation Selection Strategy}
\label{sec:repr:strategy}

Given the characteristics summarised in \Cref{tab:routeA_summary}
and~\Cref{tab:route_comparison}, the experimental campaign should follow
a principled selection strategy:

\begin{enumerate}
    \item \textbf{Start with Route~A.}  Establish baseline performance
        using the Flattened Three-Channel encoding~(A2) with classical models,
        and Three-Channel Hybrid encoding~(A1) with 1D CNNs.
        This provides a low-complexity reference against which all
        subsequent representations and algorithms are compared.
    \item \textbf{Incrementally add information.}  Evaluate A3
        (Derivative Augmented) to quantify the marginal value of local-shape features.
    \item \textbf{Introduce Route~B only if justified.}  Route~B should
        be explored if Route~A performance plateaus and there is a
        hypothesis that 2D spatial patterns (e.g.\ inter-frequency
        correlations visible in heatmaps) contain discriminative
        information that 1D convolutions cannot capture.
    \item \textbf{Always compare against Route~A.}  Any Route~B result
        must be benchmarked against the corresponding Route~A encoding to
        disentangle genuine spectral information from rendering
        artefacts.
\end{enumerate}

\noindent
This incremental strategy avoids premature commitment to high-complexity
representations and ensures that every added dimension of information
earns its place through measurable performance improvement.
